% Source: https://artofproblemsolving.com/community/c3h1218619_geometry_trivial
% Stable identity: aops-topic-1218619-q01
The four quarter discs have total area $4\cdot\dfrac{1}{4}\pi\cdot12=12\pi$. Since $2\sqrt3>3$, adjacent quarter discs overlap. Two circles of radius $r=2\sqrt3$ with centres $6$ apart meet in a lens; the half-chord satisfies $\cos\theta=\dfrac{3}{2\sqrt3}=\dfrac{\sqrt3}{2}$, so each centre sees the lens under $2\theta=60^\circ$, and
\[\text{lens}=2\left(\dfrac{60}{360}\pi r^2-\dfrac{1}{2}r^2\sin60^\circ\right)=2\left(2\pi-3\sqrt3\right)=4\pi-6\sqrt3.\]
Half of each lens lies inside the square: $2\pi-3\sqrt3$. Opposite corners are $6\sqrt2>4\sqrt3$ apart, so their discs do not meet. Covered area: $12\pi-4\left(2\pi-3\sqrt3\right)=4\pi+12\sqrt3$. Uncovered:
\[\boxed{36-4\pi-12\sqrt3}\approx2.65.\]

\begin{asy}
usepackage("amsmath");
size(3cm);
defaultpen(fontsize(10pt));
real r=2*sqrt(3); draw((0,0)--(6,0)--(6,6)--(0,6)--cycle,linewidth(1));
draw(arc((0,0),r,0,90)); draw(arc((6,0),r,90,180)); draw(arc((6,6),r,180,270)); draw(arc((0,6),r,270,360));
label("$6$",(3,0),S);
label("\mbox{}",(min(currentpicture,user=true)+max(currentpicture,user=true))/2,black); // keeps the text after the figure black when LaTeX typesets the labels inline
\end{asy}
